\documentclass[../main.tex]{subfiles}

\begin{document}
\chapter{Orientations}\label{ch:orientations}

Simply speaking, the orientation of a line or curve is a choice of direction along it, but we can generalize this idea to higher-dimensional manifolds.

\section{Orientation of Vector Spaces}

\begin{definition}{Orientation of a Vector Space}{Orientation of a Vector Space}
  Let $V$ be a $n$-dimensional real vector space. We say that two basis $(E_1, \ldots , E_n)$ and $(\tilde{E}_1, \ldots , \tilde{E}_n)$ of $V$ are \textbf{consistently oriented} if the transition matrix $B$ defined by
  \begin{equation*}
    E_i = B^j_i \tilde{E}_j
  \end{equation*}
  has positive determinant. If $n=0$, we define orientations by a choice of $+1$ or $-1$ specially.

  It is easy to see that this is an equivalence relation, and thus we can define an \textbf{orientation} of $V$ to be one of the two equivalence classes, denoted by $[E_1, \ldots , E_n]$ and its opposite $-[E_1, \ldots , E_n]$. A vector space with a choice of orientation is called an \textbf{oriented vector space}, and the chosen orientation is called \textbf{oriented} or \textbf{positively oriented}, while the opposite orientation is called \textbf{negatively oriented}.
\end{definition}

\begin{example}{Orientated Vector Spaces}{Orientated Vector Spaces}
  The orientation $[e_1, \ldots , e_n]$ of $\mathbb{R}^n$ is called the \textbf{standard orientation} of $\mathbb{R}^n$. An oriented basis for $\mathbb{R}^3$ is called a \textbf{right-handed} basis.
\end{example}

There is a connection between orientations and alternating tensors.
\begin{proposition}{Orientation and Alternating Tensors}{Orientation and Alternating Tensors}
  Let $V$ be a $n$-dimensional real vector space. For any nonzero $\omega \in \Lambda^n(V^*)$, we can define an orientation $\mathcal{O}_\omega$ of $V$ by:
  \begin{itemize}
    \item If $n\geq 1$, then $[E_1, \ldots , E_n] = \mathcal{O}_\omega$ if and only if $\omega(E_1, \ldots , E_n) > 0$.
    \item If $n=0$, then $\mathcal{O}_\omega = +1$ if $\omega > 0$ and $\mathcal{O}_\omega = -1$ if $\omega < 0$.
  \end{itemize}
\end{proposition}
\begin{proof}
  We assume $n \geq 1$, and $(E_i)$ and $(\tilde{E}_i)$ are two bases of $V$ and let $B$ be the linear map sending $E_i$ to $\tilde{E}_i$, or $\tilde{E}_j = B^i_j E_i$. Then we have
  \begin{equation*}
    \omega(\tilde{E}_1, \ldots , \tilde{E}_n) = \omega(B^i_1 E_i, \ldots , B^i_n E_i) = \det(B) \omega(E_1, \ldots , E_n).
  \end{equation*}
  So we have $\omega(\tilde{E}_1, \ldots , \tilde{E}_n) > 0$ if and only if $\det(B) > 0$ and $\omega(E_1, \ldots , E_n) > 0$, which means that the two bases are consistently oriented.
\end{proof}
If $\omega$ here determines the positive orientation, then we say that $\omega$ is a \textbf{positive orientated $n$-covector} of $V$. As $\Lambda^n(V^*)$ is one-dimensional, so choosing orientation for $V$ is equivalent to choosing a component of $\Lambda^n(V^*) \setminus \{0\}$. This formula also shows why we define the 0-dimensional case like this.

\section{Orientation of Manifolds}
\begin{definition}{Orientation of Manifolds}{Orientation of Manifolds}
  Let $M$ be a smooth manifold, with or without boundary. We define a \textbf{pointwise orientation} of $M$ to be a choice of orientation for each tangent space $T_pM$ at each point $p \in M$.

  To connect different points, we give the following definitions.
  \begin{itemize}
    \item If $(E_i)$ is a local frame for $TM$, we say that $(E_i)$ is \textbf{positively oriented} if $(E_i|_p)$ is positively oriented for each $p \in U \subseteq M$. Same for negatively oriented.
    \item A pointwise orientation of $M$ is called \textbf{continuous} if for each $p \in M$, there exists a local frame $(E_i)$ defined on a neighborhood of $p$ such that $(E_i)$ is positively oriented.
  \end{itemize}
  An \textbf{orientation} of $M$ is a continuous pointwise orientation. We say $M$ is \textbf{orientable} if it admits an orientation, and \textbf{non-orientable} otherwise. An oriented manifold is a pair $(M, \mathcal{O})$ where $M$ is a manifold and $\mathcal{O}$ is an orientation of $M$. An oriented manifold with boundary is defined similarly.
\end{definition}

Specially, if $M$ is 0-dimensional (discrete), then an orientation of $M$ is a choice of $+1$ or $-1$ for each point in $M$, and the continuity condition is vacuous. Every 0-dimensional manifold is orientable.

If $n \geq 1$, then every connected domain local frame is either positively oriented or negatively oriented.

We can also define orientation by $n$-forms.
\begin{proposition}{Orientation by $n$-forms}{Orientation by $n$-forms}
  Let $M$ be a smooth $n$-dimensional manifold, with or without boundary. For any nonvanishing $\omega \in \Omega^n(M)$, we can define an orientation $\mathcal{O}_\omega$ of $M$ by being positively oriented at each point $p$. Conversely, if $M$ is given an orientation $\mathcal{O}$, then there exists a nonvanishing $\omega \in \Omega^n(M)$ such that $\mathcal{O} = \mathcal{O}_\omega$.
\end{proposition}
\begin{proof}
  Assume $\omega$ bes a nonvanishing $n$-form on $M$, then $\omega$ defines a pointwise orientation $\mathcal{O}_\omega$ of $M$ by being positively oriented at each point. We need to show that $\mathcal{O}_\omega$ is continuous. It is trivial when $n=0$, so we assume $n \geq 1$. For each $p \in M$, we can choose a local frame $(E_i)$ defined on a connected neighborhood $U$ of $p$ and let $(\epsilon^i)$ be the dual coframe. Then we can write $\omega = f \epsilon^1 \wedge \cdots \wedge \epsilon^n$ for some smooth function $f$ on $U$. Since $\omega$ is nonvanishing, we have
  \begin{equation*}
    \omega(E_1, \ldots , E_n) = f \neq 0
  \end{equation*}
  at all points in $U$. So from connectedness of $U$, we have $f > 0$ or $f < 0$ at all points in $U$. So $(E_i)$ is positively oriented or negatively oriented at all points in $U$, which means that $\mathcal{O}_\omega$ is continuous.

  Conversely, assume $\mathcal{O}$ is an orientation of $M$. Define $\Lambda_+^n T^*M \subseteq \Lambda^n T^*M$ to be the open subset consisting of all positive orientated $n$-covectors at each point of $M$. At each $p\in M$, the intersection of $\Lambda_+^n T^*M$ with the fiber $\Lambda^n T_p^*M$ is a half line, and therefore convex. So locally in some trivialization, we can find a smooth section of $\Lambda_+^n T^*M$ over some neighborhood of $p$.

  So we can choose a smooth local section $\omega_\alpha$ of $\Lambda_+^n T^*M$ over each element $U_\alpha$ of an open cover of $M$. Then we can choose a partition of unity $\{\rho_\alpha\}$ subordinate to $\{U_\alpha\}$, and define
  \begin{equation*}
    \omega = \sum_\alpha \rho_\alpha \omega_\alpha.
  \end{equation*}
  Then $\omega$ is a smooth section of $\Lambda_+^n T^*M$, and thus a nonvanishing $n$-form on $M$. It is easy to check that $\mathcal{O} = \mathcal{O}_\omega$.
\end{proof}

A smooth coordinate chart $(U, \varphi)$ of $M$ is called \textbf{positively oriented} if the local coordinate frame $(\partial/\partial x^i)$ is positively oriented. A smooth atlas $\{(U_\alpha, \varphi_\alpha)\}$ of $M$ is called \textbf{consistently oriented} if the transition maps $\varphi_\beta \circ \varphi_\alpha^{-1}$ have positive Jacobian determinant at each point of their domains.

\begin{proposition}{Orientation by Coordinate Atlas}{Orientation by Coordinate Atlas}
  Let $M$ be a smooth manifold with positive dimension, with or without boundary. Given any consistently oriented smooth atlas $\{(U_\alpha, \varphi_\alpha)\}$ of $M$, there is a unique orientation $\mathcal{O}$ of $M$ such that each chart in the atlas is positively oriented.

  Conversely, if $M$ is given an orientation $\mathcal{O}$, and either $\partial M = \emptyset $ or $\dim M \geq 2$, then the collection of all positively oriented charts forms a consistently oriented smooth atlas of $M$.
\end{proposition}
\begin{proof}
  All is very obvious. But we shall note that for smooth 1-dimensional manifolds with boundary, the problem exists because we require the last coordinate of a boundary chart to be nonnegative, which means that if we get a negatively oriented chart at the boundary, then we cannot change it to a positively oriented one by changing the sign. For the rest cases, we can always change the sign of $x^1$ to get a positively oriented chart.
\end{proof}

\begin{proposition}{Product Orientation}{Product Orientation}
  Let $M_1, \ldots ,M_k$ be orientable smooth manifolds. There is a unique orientation on the product manifold $M_1 \times \cdots \times M_k$ such that for each $\omega_i$ that are orientation form for $M_i$, the form $\pi_1^* \omega_1 \wedge \cdots \wedge \pi_k^* \omega_k$ is an orientation form for $M_1 \times \cdots \times M_k$, where $\pi_i : M_1 \times \cdots \times M_k \to M_i$ is the projection map.
\end{proposition}
\begin{proof}
  Obvious.
\end{proof}

\begin{proposition}{Orientations of Orientable Manifolds}{Orientations of Orientable Manifolds}
  Let $M$ be a connected orientable smooth manifold, with or without boundary. Then $M$ has exactly two orientations. If two orientations $\mathcal{O}$ and $\tilde{\mathcal{O}}$ of $M$ agree at some point $p \in M$, then $\mathcal{O} = \tilde{\mathcal{O}}$.
\end{proposition}

\begin{proposition}{Orientation of Codimension 0 Submanifolds}{Orientation of Codimension 0 Submanifolds}
  Let $M$ be a smooth manifold, with or without boundary, and $N \subseteq M$ be a codimension 0 submanifold. If $M$ is orientable, then $N$ is orientable. Moreover, if $\mathcal{O}$ is an orientation of $M$, then the restriction $\mathcal{O}|_N$ is an orientation of $N$. If $\omega$ is an orientation form for $M$, then $\iota_N^* \omega$ is an orientation form for $N$, where $\iota_N : N \to M$ is the inclusion map.
\end{proposition}

Let $M,N$ be oriented smooth manifolds, with or without boundary. Let $F: M \to N$ be a local diffeomorphism. We say that $F$ is \textbf{orientation-preserving} if for each $p \in M$, the differential $dF_p : T_pM \to T_{F(p)}N$ is an orientation-preserving linear isomorphism, i.e. $dF_p$ sends positively oriented bases of $T_pM$ to positively oriented bases of $T_{F(p)}N$. The same goes for \textbf{orientation-reversing}.

\begin{proposition}{Equivalent Criterion for Orientation-Preserving Local Diffeomorphisms}{Equivalent Criterion for Orientation-Preserving Local Diffeomorphisms}
  Let $M,N$ be oriented smooth manifolds with positive dimension, with or without boundary. Let $F: M \to N$ be a local diffeomorphism. Then the following are equivalent:
  \begin{itemize}
    \item $F$ is orientation-preserving.
    \item For any oriented smooth charts for $M$ and $N$, the Jacobian of $F$ has positive determinant.
    \item For any positively oriented form $\omega$ for $N$, the pullback $F^* \omega$ is a positively oriented form for $M$.
  \end{itemize}
\end{proposition}

We can also construct orientations by pullbacks.
\begin{proposition}{The Pullback Orientation}{The Pullback Orientation}
  Suppose $M,N$ are smooth manifolds, with or without boundary, and $F: M \to N$ is a local diffeomorphism. If $N$ is oriented by $\mathcal{O}$, then $M$ has a unique orientation $F^*\mathcal{O}$ called the \textbf{pullback orientation} of $M$ by $F$ such that $F$ is orientation-preserving.
\end{proposition}
\begin{proof}
  For each $p\in M$, there is a unique orientation of $T_pM$ such that $dF_p : T_pM \to T_{F(p)}N$ is orientation-preserving. This defines a pointwise orientation of $M$. To show that it is continuous, we choose an orientation form $\omega$ for $N$. Then $F^* \omega$ is a nonvanishing $n$-form on $M$, and thus defines a continuous orientation of $M$.
\end{proof}

For this we have the usual pullback composition formula:
\begin{equation}
  (G \circ F)^* \mathcal{O} = F^*(G^* \mathcal{O})
\end{equation}

Recall that a smooth manifold $M$ is called \textbf{parallelizable} if it admits a smooth global frame for.
\begin{proposition}{Parallelizable Manifolds are Orientable}{Parallelizable Manifolds are Orientable}
  Every parallelizable smooth manifold is orientable.
\end{proposition}
\begin{proof}
  Just define the orientation by the global frame.
\end{proof}

So $\mathbb{R}^n, T^n, S^1, S^3, S^7$ and there products are orientable, and any codimension 0 submanifolds of them are orientable. Also, every Lie group is orientable because it is parallelizable. We say that an orientation of a Lie group $G$ is \textbf{left-invariant} if every left translation $L_g : G \to G$ is orientation-preserving.

\begin{proposition}{Orientations of Lie Groups}{Orientations of Lie Groups}
  Every Lie group has exactly two left-invariant orientations, corresponding to the two orientations of its Lie algebra.
\end{proposition}
\begin{proof}
  Quite natural, just define the orientation at $\Lie(G) \cong T_eG$ and then extend it to the whole $G$ by left translations.
\end{proof}

\subsection{Orientation of Hypersurfaces}
An orientation of some smooth manifold $M$ do not generally restrict to a submanifold easily. A directly restriction of a full rank form to a lower dimensional submanifold is just zero. As it turns out, submanifolds of oriented manifolds may not actually be orientable, take the M\"obius strip as an example.

Instead, here we focus on the case of hypersurfaces, which are codimension 1 submanifolds, with one extra information: A vector field that is nowhere tangent to the hypersurface.

Recall that a vector field along $S \subseteq M$ is a section of the ambient tangent bundle $TM|_S$, which is a continuous map $N : S \to TM$ such that $N_p \in T_pM$ for each $p \in S$. 

\begin{theorem}{Orientation of Hypersurfaces}{Orientation of Hypersurfaces}
  Let $M$ be an oriented smooth manifold, with or without boundary, and $S \subseteq M$ be an immersed hypersurface with or without boundary, and $N$ is a vector field along $S$ that is nowhere tangent to $S$. Then there is a unique orientation $\mathcal{O}$ of $S$ such that for each $p \in S$, a basis $(E_1, \ldots , E_{n-1})$ of $T_pS$ is positively oriented if and only if the basis $(N_p, E_1, \ldots , E_{n-1})$ of $T_pM$ is positively oriented.

  If $\omega$ is an orientation form for $M$, then $\iota_S^* (N \lrcorner \, \omega)$ is an orientation form for $S$ with respect to $\mathcal{O}$, where $\iota_S : S \hookrightarrow M$ is the inclusion map.
\end{theorem}
For $n=1$ case, since $S$ is 0-dimensional, we interpret the condition as follows: $+1$ is the positive orientation of $T_pS$ if and only if $(N_p)$ is a positively oriented basis of $T_pM$.

\begin{proof}
  We already knows that $\sigma = \iota_S^* (N \lrcorner \, \omega)$ is an $(n-1)$-form on $S$. We need to show that $\sigma$ is nonvanishing: Take any $p \in S$ and a basis $(E_1, \ldots , E_{n-1})$ of $T_pS$. Then $(N_p, E_1, \ldots , E_{n-1})$ is a basis of $T_pM$ because $N_p$ is not tangent to $S$. So we have
  \begin{equation*}
    \sigma_p(E_1, \ldots , E_{n-1}) = \omega_p(N_p, E_1, \ldots , E_{n-1}) \neq 0.
  \end{equation*}
  So $\sigma$ is nonvanishing, and thus defines an orientation $\mathcal{O}$ of $S$.
\end{proof}

However, not all hypersurfaces have a nowhere tangent vector field, however for most cases we can. For example, $S^n \subseteq \mathbb{R}^{n+1}$ has a nowhere tangent vector field given by $N = x^i \partial/\partial x^i$, which defines the standard orientation of $S^n$.

\begin{proposition}{Regular Level Sets are Orientable}{Regular Level Sets are Orientable}
  Let $M$ be an oriented smooth manifold, with or without boundary, and $f : M \to \mathbb{R}$ be a smooth function. If $S \subseteq M$ is a regular level set of $f$, then $S$ is orientable.
\end{proposition}
\begin{proof}
  Choose any Riemannian metric on $M$, and let $N = \grad f|_S$ would do.
\end{proof}

\subsection{Boundary Orientation}
If $M$ is a smooth manifold with boundary, then $\partial M$ is an embedded hypersurface of $M$. And we know from lemma \ref{lem:Existence of Inward Vector Fields}

\begin{proposition}{Induced Orientation on the Boundary}{Induced Orientation on the Boundary}
  Let $M$ be an oriented smooth $n$-dimensional manifold with boundary, $n \geq 1$. Then $\partial M$ is orientable, and all outward-pointing vector fields along $\partial M$ define the same orientation of $\partial M$, called the \textbf{induced orientation} or \textbf{Stokes orientation} of $\partial M$.
\end{proposition}

\begin{example}{Induced Orientation on the Boundary}{Induced Orientation on the Boundary}
  First, consider $S^n$ to be the boundary of the unit ball in $\mathbb{R}^{n+1}$, then the induced orientation of $S^n$ is the standard orientation.

  Next, consider the upper half space $\mathbb{H}^n$. Identify $\partial \mathbb{H}^n \subseteq \mathbb{R}^n$ by $\{(x^1, \ldots , x^{n-1}, 0) : x^i \in \mathbb{R}\}$, then the vector field $- \partial / \partial x^n$ is outward-pointing along $\partial \mathbb{H}^n$, and in $\mathbb{R}^n$, we have
  \begin{equation*}
    \left[ -\frac{\partial}{\partial x^n}, \frac{\partial}{\partial x^1}, \ldots , \frac{\partial}{\partial x^{n-1}} \right] = (-1)^n [e_1, \ldots , e_n].
  \end{equation*}
  Thus, the induced orientation of $\partial \mathbb{H}^n$ is equal to the standard orientation of $\mathbb{R}^{n-1}$ if $n$ is even, and the opposite of the standard orientation if $n$ is odd.
\end{example}

In many settings, we use parametrization to describe submanifolds.
\begin{lemma}{Local Parametrization of Boundary and Orientation}{Local Parametrization of Boundary and Orientation}
  Let $M$ be an oriented smooth $n$-dimensional manifold with boundary. Suppose $U \subseteq \mathbb{R}^{n-1}$ is an open subset, and $a,b\in \mathbb{R}, a<b$, and $F: (a,b] \times U \to M$ is a smooth embedding that restricts to a smooth embedding of $\{b\} \times U$ into $\partial M$. Then the parametrization $f: U \rightarrow \partial M, f(x) = F(b,x)$ is orientation-preserving for $\partial M$ if and only if $F$ is orientation-preserving for $M$.
\end{lemma}
\begin{proof}
  Quite natural, note that a same-dimension embedding is a diffeomorphism onto its image, so we can identify $U$ with $\{b\} \times U$ and then the condition is just the same as the definition of orientation-preserving for $F$.
\end{proof}

For example, consider the smooth local parametrization of $S^2 = \partial \overline{B}^3$, let $U = (0, \pi) \times (0, 2\pi)$ and $X: U \to S^2 \subseteq \mathbb{R}^3$ be given by
\begin{equation*}
  X(\theta, \varphi) = (\sin \theta \cos \varphi, \sin \theta \sin \varphi, \cos \theta).
\end{equation*}
Then it is a restriction of a smooth embedding $F: (0, 1] \times U \to \overline{B}^3$ given by
\begin{equation*}
  F(\rho, \theta, \varphi) = (\rho \sin \theta \cos \varphi, \rho \sin \theta \sin \varphi, \rho \cos \theta), \qquad F(1, \theta, \varphi) = X(\theta, \varphi).
\end{equation*}
So, as $\det J_F = \rho^2 \sin \theta > 0$, we have $X$ is orientation-preserving for $S^2$.

\section{The Riemannian Volume Form}
Let $(M,g)$ be an oriented Riemannian manifold of positive dimension. Then there is a smooth orthonormal local frame $(E_1, \ldots , E_n)$ in a neighborhood of each point in $M$. Assume it is positively oriented (else we can change the sign of $E_1$ to make it positively oriented), then we can find an oriented orthonormal frame in a neighborhood of each point in $M$.

\begin{proposition}{Riemannian Volume Form}{Riemannian Volume Form}
  Suppose $(M,g)$ is an oriented Riemannian $n$-dimensional manifold, with or without boundary, and $n \geq 1$. Then there is a unique orientation form $\omega_g \subseteq \Omega^n(M)$ called the \textbf{Riemannian volume form}, such that
  \begin{equation}
    \omega_g(E_1, \ldots , E_n) = 1.
  \end{equation}
  for every local oriented orthonormal frame $(E_1, \ldots , E_n)$ for $M$.

  In local smooth oriented coordinates $(U, x^i)$, we have
  \begin{equation}
    \omega_g = \sqrt{\det(g_{ij})} dx^1 \wedge \cdots \wedge dx^n,
  \end{equation}
\end{proposition}
\begin{proof}
  Existence and uniqueness is easy. To see the coordinate form, let $(E_i)$ be the local oriented orthonormal frame, then let
  \begin{equation*}
    \frac{\partial}{\partial x^i} = A_i^j E_j, \qquad \omega_g = f dx^1 \wedge \cdots \wedge dx^n.
  \end{equation*}
  Then we have
  \begin{equation*}
    f = \omega_g\left( \frac{\partial}{\partial x^1}, \ldots , \frac{\partial}{\partial x^n} \right) = \omega_g(A_1^j E_j, \ldots , A_n^j E_j) = \det(A).
  \end{equation*}
  \begin{equation*}
    g_{ij} = g\left( \frac{\partial}{\partial x^i}, \frac{\partial}{\partial x^j} \right) = g(A_i^k E_k, A_j^l E_l) = A_i^k A_j^l g(E_k, E_l) = \sum_{k=1}^n A_i^k A_j^k,
  \end{equation*}
  So $\det(g_{ij}) = \det(A)^2$, and thus $f = \sqrt{\det(g_{ij})}$, since both frames are positively oriented. (If $(x^i)$ is a negatively oriented coordinate chart, then we have $f = -\sqrt{\det(g_{ij})}$, which is consistent with the definition of $\omega_g$.)
\end{proof}

This form is also isometric invariant. Suppose $(M, g)$ and $(\tilde{M}, \tilde{g})$ are positive-dimensional oriented Riemannian manifolds, with or without boundary, and $F: M \to \tilde{M}$ is a local isometry, then $F^* \omega_{\tilde{g}} = \omega_g$.

\subsection{Hypersurfaces in Riemannian Manifolds}
Let $(M,g)$ be an oriented Riemannian manifold, with or without boundary, and $S \subseteq M$ be an immersed hypersurface with or without boundary. Any unit normal vector field $N$ along $S$ is nowhere tangent to $S$, so it defines an orientation of $S$ by previous discussions.

\begin{proposition}{Volume Form on Submanifolds}{Volume Form on Submanifolds}
  Let $(M,g)$ be an oriented Riemannian manifold, with or without boundary, and $S \subseteq M$ be an immersed hypersurface with or without boundary. Let $\tilde{g}$ be the induced Riemannian metric on $S$, and suppose $N$ is a unit normal vector field along $S$. Then with respect to the orientation of $S$ defined by $N$, the Riemannian volume form of $(S, \tilde{g})$ is given by
  \begin{equation}
    \omega_{\tilde{g}} = \iota_S^* (N \lrcorner \, \omega_g),
  \end{equation}
\end{proposition}

\begin{proposition}{Outward Pointing Normal Field}{Outward Pointing Normal Field}
  Let $(M,g)$ be a Riemannian manifold with boundary. There is a unique smooth outward-pointing unit normal vector field $N$ along $\partial M$.
\end{proposition}
\begin{proof}
  Existence: Let $f$ be a boundary defining function for $M$, then
  \begin{equation*}
    N = -\frac{\grad f}{|\grad f|_g} \bigg|_{\partial M}
  \end{equation*}
  is a smooth outward-pointing unit normal vector field along $\partial M$. Uniqueness is obvious.
\end{proof}

\begin{corollary}{Structure of the Boundary Volume Form}{Structure of the Boundary Volume Form}
  Let $(M,g)$ be an oriented Riemannian manifold with boundary, and $\tilde{g}$ be the induced Riemannian metric on $\partial M$. Then the Riemannian volume form of $(\partial M, \tilde{g})$ is
  \begin{equation*}
    \omega_{\tilde{g}} = \iota_{\partial M}^* (N \lrcorner \, \omega_g),
  \end{equation*}
  where $N$ is the unique smooth outward-pointing unit normal vector field along $\partial M$.
\end{corollary}

\section{Orientations and Covering Maps}
Proving that a manifold is non-orientable is usually more difficult than proving that it is orientable. One useful tool for this is the following result about covering maps.

First here is a simply special case of the pullback orientation:
\begin{proposition}{Orientation of Covering Spaces}{Orientation of Covering Spaces}
  If $\pi: E \to M$ is a smooth covering map, and $M$ is orientable, then $E$ is orientable.
\end{proposition}
\begin{proof}
  Covering maps are local diffeomorphisms, so it follows from proposition \ref{prop:The Pullback Orientation}.
\end{proof}

If $G$ is a Lie group acting smoothly on a smooth manifold $E$, we say that the action is \textbf{orientation-preserving} if for each $g \in G$, the diffeomorphism $E \to E, p \mapsto g \cdot p$ is orientation-preserving.

A covering map $\pi: E \rightarrow M$ is called \textbf{normal} (also \textbf{regular} or \textbf{Galois}) if the group $G$ of deck transformations of $\pi$ acts transitively on each fiber of $\pi$. In this case, we have $M \cong E/G$. Here $G$ is denoted by $\Aut_\pi(E)$, and is called the \textbf{automorphism group} of $\pi$.

\begin{theorem}{Orientability and Normal Covering Maps}{Orientability and Normal Covering Maps}
  Suppose $E$ is a connected, oriented smooth manifold, with or without boundary, and $\pi: E \to M$ is a normal smooth covering map. Then $M$ is orientable if and only if the action of $\Aut_\pi(E)$ on $E$ is orientation-preserving.
\end{theorem}
\begin{proof}
  Let $\mathcal{O}_E$ be the orientation of $E$. First suppose $M$ is orientable, and $q\in E$. Because $M$ is connected, it has exactly two orientations, and one of them has the property $\mathrm{d} \pi_q : T_qE \to T_{\pi(q)}M$ is orientation-preserving (because local diffeomorphisms are orientation-preserving for some orientation). Denote that orientation by $\mathcal{O}_M$. Then $\pi^* \mathcal{O}_M = \mathcal{O}_E$ by the uniqueness of pullback orientation, so for $\varphi \in \Aut_\pi(E)$, we have $\pi \circ \varphi = \pi$, and thus
  \begin{equation*}
    \varphi^* \mathcal{O}_E = \varphi^* (\pi^* \mathcal{O}_M) = (\pi \circ \varphi)^* \mathcal{O}_M = \pi^* \mathcal{O}_M = \mathcal{O}_E,
  \end{equation*}
  which means that $\varphi$ is orientation-preserving.

  Conversely, suppose the action of $\Aut_\pi(E)$ on $E$ is orientation-preserving, and let $p\in M$. If $U \subseteq M$ is any evenly covered neighborhood of $p$, then there is a smooth section $\sigma: U \to E$ that induces an orientation $\sigma^* \mathcal{O}_E$ of $U$. Suppose $\sigma': U \to E$ is another smooth section, then there is a covering automorphism $\varphi \in \Aut_\pi(E)$ such that $\sigma' = \varphi \circ \sigma$. So we have
  \begin{equation*}
    \sigma'^* \mathcal{O}_E = \sigma^* (\varphi^* \mathcal{O}_E) = \sigma^* \mathcal{O}_E,
  \end{equation*}
  because $\varphi$ is orientation-preserving. So the orientation induced by any local section is the same, and thus defines a continuous orientation of $M$ via pullback, which means that $M$ is orientable.
\end{proof}

\begin{remark}
  This is how we say that the manifold ``twists'' in the M\"obius strip, because the nontrivial deck transformation (like walking along the strip and come back to the same point) is orientation-reversing (getting to the back of the strip).
\end{remark}

\begin{example}{Orientation by Covering Maps}{Orientation by Covering Maps}
  \begin{itemize}
    \item Consider the real projective space $\mathbb{R}P^n = S^n/\{\pm 1\}$, where the only nontrivial covering automorphism is given by $\alpha(x) = -x$. We will see that $\alpha$ is orientation-preserving if and only if $n$ is odd, so $\mathbb{R}P^n$ is orientable if and only if $n$ is odd.
    \begin{proof}
      Consider $F: \overline{B}^{n+1} \to \overline{B}^{n+1}$ given by $F(x) = -x$, and take the Riemannian volume form $\omega_g$ of $\overline{B}^{n+1}$, then the induced orientation of $S^n$ is the standard orientation by $\omega_{\tilde{g}} = \iota_{S^n}^* (N \lrcorner \, \omega_g)$, where $N$ is the outward-pointing normal vector field along $S^n$. Consider the pullback $\alpha^* \omega_{\tilde{g}}$, we have $\iota_{S^n} \circ \alpha = F \circ \iota_{S^n}$, so
      \begin{equation*}
        \begin{aligned}
          \alpha^* \omega_{\tilde{g}} & = \alpha^* (\iota_{S^n}^* (N \lrcorner \, \omega_g)) = (\iota_{S^n} \circ \alpha)^* (N \lrcorner \, \omega_g)\\
          & = (F \circ \iota_{S^n})^* (N \lrcorner \, \omega_g) = \iota_{S^n}^* (F^* (N \lrcorner \, \omega_g)) = \iota_{S^n}^* ((dF^{-1} N) \lrcorner \, F^* \omega_g)\\
          & = \iota_{S^n}^* (N \lrcorner \, (-1)^{n+1} \omega_g) = (-1)^{n+1} \iota_{S^n}^* (N \lrcorner \, \omega_g) = (-1)^{n+1} \omega_{\tilde{g}}.
        \end{aligned}
      \end{equation*}
      where we used
      \begin{equation*}
        (\mathrm{d} F^{-1} N)_x(f) = N_{F^{-1}(x)}(f \circ F^{-1}) = N_{-x}(f \circ F^{-1}) = N_x(f)
      \end{equation*}
      so $\alpha^* \omega_{\tilde{g}} = \omega_{\tilde{g}}$ if $n$ is odd, and $\alpha^* \omega_{\tilde{g}} = -\omega_{\tilde{g}}$ if $n$ is even.
    \end{proof}

  \item Now consider the M\"obius band. Let $E$ be the total space of M\"obius bundle. The quotient map $q: \mathbb{R}^2 \to E$. Then $\Aut_q(\mathbb{R}^2) \cong \mathbb{Z}$ by $n \cdot (x,y) = (x + n, (-1)^n y)$. So for odd $n$, the pullback gives $ \mathrm{d} x \wedge \mathrm{d} y \mapsto - \mathrm{d} x \wedge \mathrm{d} y$, which is orientation-reversing, so $E$ is non-orientable.
  \end{itemize}
\end{example}

\subsection{The Orientation Covering}
Now we shall see that every non-orientable manifold has a two-sheeted covering manifold. This gives intuition that every non-orientable manifold is like a twisted version of an orientable manifold.

Recall that a topological covering map $\pi: E \to M$ between \textbf{connected and locally path connected topological spaces} is defined as a surjective continuous map such that for each $p \in M$, there is an open neighborhood $U$ of $p$ such that $\pi^{-1}(U)$ is a disjoint union of open subsets of $E$, each of which is homeomorphic to $U$ by $\pi$. A smooth covering map is a topological covering map that is also a smooth map.

We define a \textbf{generalized covering map} to be covering map except that we do not require the total space $E$ to be connected. Well, we can easily see that its restriction to each connected component of $E$ is a covering map in the usual sense.

\begin{definition}{Orientation Covering}{Orientation Covering}
  Let $M$ be a connected smooth positive-dimensional manifold, with or without boundary. Let $\hat{M}$ be the set of all orientations of all tangent spaces of $M$, i.e.
  \begin{equation*}
    \hat{M} = \{(p, \mathcal{O}_p) : p \in M, \mathcal{O}_p \text{ is an orientation of } T_pM\}.
  \end{equation*}
  Define the projection map $\hat{\pi} : \hat{M} \to M$ by $\hat{\pi}(p, \mathcal{O}_p) = p$. Then $\hat{\pi}$ is called the \textbf{orientation covering} of $M$, each fiber $\hat{\pi}^{-1}(p)$ consists of exactly two points.
\end{definition}

\begin{proposition}{Properties of the Orientation Covering}{Properties of the Orientation Covering}
  Let $M$ be a connected smooth positive-dimensional manifold, with or without boundary, and $\hat{\pi} : \hat{M} \to M$ be the orientation covering of $M$. Then $\hat{M}$ can be given s structure of smooth oriented manifold with or without boundary such that
  \begin{itemize}
    \item $\hat{\pi}$ is a generalized smooth covering map.
    \item A connected open subset $U \subseteq M$ is evenly covered by $\hat{\pi}$ if and only if $U$ is orientable. In this case, every orientation of $U$ is the pullback orientation induced by a local section of $\hat{\pi}$ over $U$.
  \end{itemize}
\end{proposition}

\begin{remark}
  We can view $\hat{M}$ as the two sides of $M$, geometrically, imaging a smooth surface, and moving its two sides apart, then the resulting manifold is just $\hat{M}$, and the projection map $\hat{\pi}$ just glues the two sides together. For orientable manifolds like $S^2$, the two sides are completely separated, so $\hat{M}$ is just two disjoint copies of $M$. For non-orientable manifolds like the M\"obius strip, the two sides are twisted together, but $\hat{M}$ is actually orientable because it needs two rounds to get back to the same point.
\end{remark}
\begin{proof}
  \begin{itemize}
    \item First, we give $\hat{M}$ a topology: For each orientable open subset $U \subseteq M$ and each orientation $\mathcal{O}$ of $U$, we define
      \begin{equation*}
        \hat{U}_\mathcal{O} = \{(p, \mathcal{O}_p) : p \in U, \mathcal{O}_p = \mathcal{O}|_p\} \subseteq \hat{M}.
      \end{equation*}
      Then all $\hat{U}_\mathcal{O}$ forms a basis of a topology on $\hat{M}$, and $\hat{\pi}|_{\hat{U}_\mathcal{O}} : \hat{U}_\mathcal{O} \to U$ is a homeomorphism, so $\hat{\pi}$ is a local homeomorphism. Now, as every orientable connected open subset $U \subseteq M$ has $\hat{\pi}^{-1}(U) = \hat{U}_\mathcal{O} \sqcup \hat{U}_{-\mathcal{O}}$, so $\hat{\pi}$ is a generalized covering map. By the preceding discussion, $\hat{\pi}$ restricted to each connected component of $\hat{M}$ is a covering map in the usual sense, so there is a unique smooth structure on $\hat{M}$ such that $\hat{\pi}$ is a smooth covering map, from theorem \ref{thm:Covering Space of Smooth manifolds}.

    \item To give $M$ an orientation, for $\hat{p} = (p, \mathcal{O}_p) \in \hat{M}$, we define $\mathcal{O}_{\hat{p}} = \mathcal{O}_p$. Then it gives a continuous orientation $\mathcal{O}$ of $\hat{M}$.
  \end{itemize}
\end{proof}

\begin{theorem}{Orientation Covering Theorem}{Orientation Covering Theorem}
  Suppose $M$ is a connected smooth manifold, with or without boundary, and $\hat{\pi} : \hat{M} \to M$ is the orientation covering of $M$.
  \begin{itemize}
    \item If $M$ is orientable, then $\hat{M}$ has exactly two connected components, each of which is diffeomorphic to $M$ by $\hat{\pi}$.
    \item If $M$ is non-orientable, then $\hat{M}$ is connected, and $\hat{\pi} : \hat{M} \to M$ is a two-sheeted smooth covering map.
  \end{itemize}
\end{theorem}

\begin{theorem}{Uniqueness of the Orientation Covering}{Uniqueness of the Orientation Covering}
  Let $M$ be a non-orientable connected smooth manifold, with or without boundary and $\pi: \hat{M} \to M$ is its orientation covering. If $\tilde{M}$ is an oriented smooth manifold, with or without boundary that exists a two-sheeted smooth covering map $\tilde{\pi} : \tilde{M} \to M$, then there is a unique orientation-preserving diffeomorphism $\varphi : \tilde{M} \to \hat{M}$ such that $\pi \circ \varphi = \tilde{\pi}$.
\end{theorem}

\begin{theorem}{Orientatability and Fundamental Group}{Orientability and Fundamental Group}
  Let $M$ be a connected smooth manifold, with or without boundary, and the fundamental group of $M$ has no subgroup of index 2, then $M$ is orientable.

  In particular, if $M$ is simply connected, then $M$ is orientable.  
\end{theorem}


\end{document}
